\section{Ausgangslage}
\setauthor{\firstauthor}
% \textit{Beschreibung der aktuellen Situation: Warum besteht Bedarf für die Matura-Lernplattform? Probleme beim Lernen, fehlendes Feedback, etc.}

Die Vorbereitung auf Tests, Schularbeiten und schließlich die Matura erfordern ein hohes Maß an Eigenverantwortung und Organisation. Im Moment erfolgt das Lernen größtenteils mithilfe von Skripten, Arbeitsblättern und Übungsbeispielen aus dem Unterricht. Diese bewährten Methoden stoßen allerdings dort an ihre Grenzen, wo es um die Dynamik des Lernprozess geht.

Zwar ist eine erfolgreiche Vorbereitung mit den vorhandenen Mitteln durchaus möglich, sie ist allerdings häufig mit einem hohen organisatorischen Aufwand verbunden. Wesentlich dabei ist das fehlende unmittelbare Feedback beim eigenständigen Üben. Schüler und Schülerinnen können ihre Ergebnisse mit Musterlösungen abgleichen, dabei bleibt jedoch der Weg zur Lösung und die Analyse spezifischer Fehlerquellen oft ungeklärt. Weiters ist es ohne digitale Unterstützung schwierig, den individuellen Lernfortschritt über einen längeren Zeitraum präzise zu verfolgen oder gezielt Schwachstellen zu identifizieren, die mehr Übung erfordern.


\subsection{Die HTL-Matura in Österreich}

Mit der Reife- und Diplomprüfung wird die Ausbildung an einer Höheren Technischen Lehranstalt (HTL) in Österreich abgeschlossen \cite{MaturaAustria}. Diese ist modular aufgebaut und in ein 3-Säulen-Modell (\autoref{fig:3saeulen}) unterteilt. Während dabei die erste Säule, die Diplomarbeit, ein in einem Team fertiggestelltes Projekt darstellt, konzentriert sich die Vorbereitung der Schüler:innen vor allem auf die weiteren zwei Säulen, die den Kern der klassischen Prüfungssituation bilden.

\begin{figure}[H]
	\centering
	\fbox{
		\includegraphics[width=0.8\textwidth]{3saeulen.jpg}
	}
	\caption{Das 3-Säulen-Modell der Reife- und Diplomprüfung (\cite{MaturaAustria})}
	\label{fig:3saeulen}
\end{figure}


\textbf{Die schriftliche Klausurprüfung (Säule 2)}

Im schriftlichen Teil der Matura müssen Schüler:innen entweder drei oder vier Klausurarbeiten absolvieren, je nachdem, ob sie alle vier Fächer schriftlich oder eine der Sprachen mündlich ablegen möchten. Hierbei gibt es zwei verschiedene Arten von Aufgaben:

\begin{itemize}
	\item \textbf{Zentralisierte Standardisierung:} \\
	In den Hauptfächern Deutsch, Englisch und Angewandte Mathematik werden die Aufgabenstellungen vom Bundesministerium zentral vorgegeben. Das bedeutet, dass für diese Fächer in jeder Berufsbildenden Höheren Schule (BHS), wie es auch bei der HTL der Fall ist, die selben Aufgaben pro Jahrgang zur Matura kommen.
	\item \textbf{Fachspezifische Klausur:} \\
	Zusätzlich zu den Hauptfächern muss eine schriftliche Prüfung in einem schwerpunktspezifischen Fach wie beispielsweise Softwareentwicklung abgelegt werden. Diese Aufgaben erstellen die Lehrkräfte an den jeweiligen Schulstandorten.  Das heißt also, dass diese Aufgaben schulspezifisch erstellt werden und sich auch bei Schulen mit selbem Schwerpunkt unterscheiden können.
\end{itemize}


\vspace{0.4cm}
\textbf{Die mündliche Prüfung (Säule 3)}

Bei der mündlichen Säule gibt es zwei oder drei Teilprüfungen, je nachdem für welche Variante sich die Schüler:innen entschieden haben. Hierbei müssen sie aus einem Pool an Themenbereichen ziehen, die im Vorfeld angekündigt wurden. Vor allem an technischen Schulen ist dabei vorgeschrieben, dass mindestens eine Prüfung aus der Fachtheorie stammen muss. 


\vspace{0.4cm}
Bereits bestehende Lernplattformen fokussieren sich oft nur auf die zentralisierten Fächer und erfüllen die spezifischen Anforderungen des technischen Lehrplans nur bedingt. Die geplante Matura-Lernplattform setzt genau hier an. Sie soll nicht als Ersatz für den Unterricht, sondern als ergänzendes Werkzeug dienen und den Lernvorgang für Schüler:innen durch direktes Feedback und eine strukturierte Fortschrittskontrolle effizienter und transparenter gestalten.